% from ega4-15-fr.tex
\oldpage[IV-3]{241}
En outre, ces conditions entraînent que $X^0$ est ouvert dans $X$. Enfin, lorsque $Y$ est localement noethérien, a) et b) sont aussi équivalentes à
\begin{enumerate}
  \item[\rm b$'$)] $f$ est universellement ouvert aux points de $X^0$.
\end{enumerate}
\end{corollary}

\begin{proof}
Le fait que a) entraîne b) résulte de (15.6.6, (ii)), ainsi que le fait que $X^0$ est alors ouvert dans $X$. Si b) est vérifiée, on peut se borner au cas où $Y$ est réduit et $f$ plat~; alors, on se ramène, comme au début de (15.6.6), et en utilisant en outre (11.2.7), au cas où $Y$ est noethérien, cas où la conclusion résulte de (15.6.6, (iii), cas $\alpha$)). L'équivalence de b) et b$'$) a déjà été prouvée lorsque $Y$ est localement noethérien (15.2.2.1), compte tenu de ce que le morphisme $X\times_Y Y_{\mathrm{red}}\to X$ est un homéomorphisme universel.
\end{proof}

\begin{proposition}[15.6.8]
\label{IV.15.6.8-fr}
Soient $f:X\to Y$ un morphisme séparé de présentation finie, $g$ une $Y$-section de $X$~; pour tout $y\in Y$, soit $X_y^0$ la composante connexe de $f^{-1}(y)$ qui contient $g(y)$, et soit $X^0$ la réunion des $X_y^0$ quand $y$ parcourt $Y$. Alors, si $y\in Y$ est tel que $X_y^0$ soit propre sur $\mathbf{k}(y)$, $X^0$ est propre sur $Y$ au point $y$ (i.e. (15.7.1), il existe un voisinage ouvert $U$ de $y$ dans $Y$ tel que $X^0\cap f^{-1}(U)$ soit fermé dans $f^{-1}(U)$ et propre sur $U$).
\end{proposition}

\begin{proof}
Procédant comme au début de (15.6.6) (où on utilise (8.10.5, (v)), et où on remplace (9.7.11) par (9.6.7)), on se ramène au cas où $Y$ est noethérien et $f$ de type fini. On sait alors (9.7.12) que $X^0$ est localement constructible dans $X$~; d'autre part, les $X_z^0$ sont géométriquement connexes (pour tout $z\in Y$) en vertu de (4.5.13)~; enfin, comme $g(Y)\subset X^0$, $X^0$ est universellement submersif sur $Y$ (15.7.8). Il suffit alors d'appliquer à $X^0$ le critère (15.7.9).
\end{proof}

\begin{corollary}[15.6.9]
\label{IV.15.6.9-fr}
Soit $f:X\to Y$ un morphisme séparé tel que $Y$ soit localement noethérien et $f$ localement de type fini, ou que $f$ soit de présentation finie. Soit $g$ une $Y$-section de $X$, et pour tout $y\in Y$, soit $X_y^0$ la composante connexe de $f^{-1}(y)$ contenant $g(y)$. Soit $y$ un point de $Y$ tel que $X_y^0$ soit propre sur $\mathbf{k}(y)$. Alors, pour toute générisation $y'$ de $y$, on a $\dim(X_y^0)\geqslant\dim(X_{y'}^0)$.
\end{corollary}

\begin{proof}
Supposons d'abord $f$ de présentation finie~; alors en remplaçant au besoin $Y$ par un voisinage de $y$, on peut supposer que $f|X^0$ est \emph{propre} (15.6.8), et il suffit d'appliquer (13.1.5) à cette restriction. Si $Y$ est localement noethérien et $f$ localement de type fini, $X$ est localement noethérien~; en outre l'hypothèse que $X_y^0$ soit propre sur $\mathbf{k}(y)$ entraîne que $X_y^0$ est noethérien~; il existe donc un ouvert noethérien $V\subset X$ contenant $X_y^0$~; par suite il y a un voisinage ouvert $W$ de $y$ dans $Y$ tel que $g(W)\subset V$. En outre, si $z$ est point maximal de $X_{y'}^0$, on peut supposer que $z\in V$. La restriction de $f$ à $V\cap f^{-1}(W)$ est alors un morphisme de type fini, et $g|W$ une $W$-section~; en appliquant à ce morphisme le résultat déjà obtenu, on voit que si $Z$ est la composante irréductible de $X_{y'}^0$ de point générique $z$, on a $\dim(Z)\leqslant\dim(X_y^0)$. Comme $z$ est un point maximal quelconque de $X_{y'}^0$, on a bien $\dim(X_{y'}^0)\leqslant\dim(X_y^0)$.
\end{proof}

\begin{remarks}[15.6.10]
\EGAReaderTarget{ega4.c04.item.15-6-10}{IV.15.6.10-fr}
(i) L'hypothèse supplémentaire sur l'existence d'une composante irréductible de $X_{y'}^0$ contenant $g(y')$ et de dimension égale à $\dim(X_{y'}^0)$, faite dans (15.6.1), n'est pas superflue, comme le montre l'exemple suivant. Soient $A$ un anneau de valuation discrète, $\pi$ une uniformisante de $A$, $K$ le corps des fractions de $A$, $k$ son corps résiduel. Prenons $Y=\operatorname{Spec}(A)$, et désignons par $y$ et $y'$ le point fermé et le point générique de $Y$. Considérons les deux $Y$-schémas suivants~: $X_1=\operatorname{Spec}(A[t])$, $X_2=\operatorname{Spec}(K[u,v])$, où $t,u,v$ sont des indéterminées (on notera que la projection de $X_2$ dans $Y$ est donc $\{y'\}$). Dans l'anneau $A[t]$, l'idéal principal $(\pi t-1)$ est \emph{maximal}, et correspond donc à un point fermé $x_1$ de $X_1$, qui se projette au point générique $y'$ de $Y$ et est tel que $\mathbf{k}(x_1)=K$.

